% =====================================================================
%  Preamble desain Modul Praktikum Pemrograman Statistika
%  Politeknik Statistika STIS - Palet warna BPS (biru, oranye, hijau)
%  Mesin: XeLaTeX
% =====================================================================

% ---------- Font ----------
\usepackage{fontspec}
% ---------- Pemenggalan kata bahasa Indonesia ----------
\language=\csname l@indonesian\endcsname

\setmainfont{Carlito}[Ligatures=TeX]
\setsansfont{Carlito}[Ligatures=TeX]
\setmonofont{DejaVu Sans Mono}[Scale=0.76]

% ---------- Paket umum ----------
\usepackage{xcolor}
\usepackage{tikz}
\usepackage{graphicx}
\usetikzlibrary{calc}
\usepackage{etoolbox}
\usepackage[most]{tcolorbox}
\usepackage{titlesec}
\usepackage{fancyhdr}
\usepackage{mdframed}
\usepackage{enumitem}
\usepackage{amsmath,amssymb}
\usepackage{microtype}
\usepackage{tocloft}
\usepackage{fancyvrb}
\usepackage{setspace}
\usepackage{ragged2e}

% ---------- Dependensi tabel (knitr/kableExtra) ----------
\usepackage{booktabs}
\usepackage{longtable}
\usepackage{array}
\usepackage{multirow}
\usepackage{wrapfig}
\usepackage{float}
\usepackage{colortbl}
\usepackage{pdflscape}
\usepackage{tabu}
\usepackage{threeparttable}
\usepackage{threeparttablex}
\usepackage[normalem]{ulem}
\usepackage{makecell}

% ---------- Palet warna BPS ----------
\definecolor{bpsnavy}{HTML}{123A6B}
\definecolor{bpsbiru}{HTML}{0B6FB4}
\definecolor{bpsbirumuda}{HTML}{6FB3E0}
\definecolor{bpsoranye}{HTML}{F58220}
\definecolor{bpshijau}{HTML}{00A651}
\definecolor{abuteks}{HTML}{45505C}
\definecolor{codebg}{HTML}{F4F7FB}
\definecolor{codebgpy}{HTML}{FFF8E8}
\definecolor{pykuning}{HTML}{E8A200}
\definecolor{outbg}{HTML}{F6FBF7}

\onehalfspacing
\setlength{\parindent}{0pt}
\setlength{\parskip}{5pt}
\setlength{\headheight}{24pt}
\color{abuteks!95!black}

% ---------- Istilah bahasa Indonesia ----------
\renewcommand{\contentsname}{Daftar Isi}
\renewcommand{\listfigurename}{Daftar Gambar}
\renewcommand{\listtablename}{Daftar Tabel}
\renewcommand{\tablename}{Tabel}
\renewcommand{\figurename}{Gambar}

% ---------- Penomoran mengikuti nomor modul ----------
\newcommand{\nomormodul}{1}
\renewcommand{\thesection}{\arabic{section}}
\renewcommand{\thesubsection}{\thesection.\arabic{subsection}}
\renewcommand{\theequation}{\arabic{equation}}
\renewcommand{\thetable}{\arabic{table}}
\renewcommand{\thefigure}{\arabic{figure}}

% ---------- Gaya judul ----------
\titleformat{\section}[hang]
  {\sffamily\bfseries\Large\color{bpsnavy}}{Bab \thesection}{0.7em}{}
  [{\vspace{1pt}{\color{bpsoranye}\titlerule[1.6pt]}}]
\titlespacing*{\section}{0pt}{18pt}{8pt}

\titleformat{\subsection}[hang]
  {\sffamily\bfseries\large\color{bpsbiru}}{\thesubsection}{0.5em}{}
\titlespacing*{\subsection}{0pt}{13pt}{4pt}

\titleformat{\subsubsection}[hang]
  {\sffamily\bfseries\normalsize\color{bpsnavy}}{\thesubsubsection}{0.5em}{}
\titlespacing*{\subsubsection}{0pt}{10pt}{3pt}

% ---------- Header & footer ----------
\pagestyle{fancy}
\fancyhf{}
\fancyhead[L]{\small\color{bpsnavy}\textbf{Pedoman \texttt{modulstis}}}
\fancyhead[R]{\small\color{abuteks}Lampiran}
\renewcommand{\headrulewidth}{1pt}
\renewcommand{\headrule}{\hbox to\headwidth{\color{bpsbiru}\leaders\hrule height \headrulewidth\hfill}}
\fancyfoot[L]{\scriptsize\color{abuteks}Politeknik Statistika STIS}
\fancyfoot[C]{\small\colorbox{bpsnavy}{\textcolor{white}{\,\thepage\,}}}
\fancyfoot[R]{\scriptsize\color{abuteks}D-IV Komputasi Statistik}
\renewcommand{\footrulewidth}{0pt}

% ---------- Daftar isi ----------
\renewcommand{\cfttoctitlefont}{\sffamily\bfseries\Large\color{bpsnavy}}
\renewcommand{\cftsecfont}{\bfseries\color{bpsnavy}}
\renewcommand{\cftsecpagefont}{\bfseries\color{bpsnavy}}
\renewcommand{\cftsubsecfont}{\color{abuteks}}
\setlength{\cftbeforesecskip}{5pt}

% ---------- Kotak-kotak ----------
\newtcolorbox{kotakinfo}[1][]{
  breakable, enhanced, colback=bpsbiru!5, colframe=bpsbiru,
  boxrule=0pt, leftrule=4pt, arc=2pt, left=10pt, right=10pt, top=8pt, bottom=8pt,
  before upper={\ifstrempty{#1}{}{\textcolor{bpsbiru!80!black}{\sffamily\bfseries #1.}~}}}

\newtcolorbox{kotaktujuan}{
  breakable, enhanced, colback=bpsnavy!5, colframe=bpsnavy,
  boxrule=0pt, leftrule=4pt, arc=2pt, left=10pt, right=10pt, top=8pt, bottom=8pt}

% Kotak penjelasan kode: menerangkan APA yang dikerjakan kode, sebelum interpretasi hasil
% Kotak anatomi: memperagakan pasangan ui-server dengan warna serasi
\definecolor{padanA}{HTML}{C2410C}   % oranye tua  - identitas MASUKAN
\definecolor{padanB}{HTML}{1D4ED8}   % biru tua    - identitas KELUARAN
\definecolor{padanC}{HTML}{047857}   % hijau tua   - objek reaktif / data
\newtcolorbox{kotakanatomi}[1][Anatomi Pasangan ui dan server]{
  breakable, enhanced, colback=white, colframe=bpsnavy, boxrule=0.7pt,
  arc=2pt, left=8pt, right=8pt, top=7pt, bottom=7pt,
  title={\sffamily\bfseries #1}, coltitle=white, colbacktitle=bpsnavy}

\newtcolorbox{kotakpenjelasan}[1][Penjelasan Kode]{
  breakable, enhanced, colback=bpsbirumuda!5, colframe=bpsnavy, boxrule=0pt, leftrule=4pt,
  arc=2pt, left=10pt, right=10pt, top=7pt, bottom=7pt,
  before upper={\textcolor{bpsnavy}{\sffamily\bfseries #1.}~}}

\newtcolorbox{kotakinterpretasi}[1][Interpretasi Output]{
  breakable, enhanced, colback=bpsbiru!6, colframe=bpsbiru, boxrule=0.9pt, arc=3pt,
  left=10pt, right=10pt, top=9pt, bottom=8pt,
  title={\sffamily\bfseries #1}, coltitle=white, colbacktitle=bpsbiru,
  fonttitle=\small, attach boxed title to top left={xshift=8pt, yshift=-3pt},
  boxed title style={boxrule=0pt, arc=2pt}}

\newtcolorbox{kotakcatatan}[1][Catatan]{
  breakable, enhanced, colback=bpshijau!7, colframe=bpshijau, boxrule=0pt, leftrule=4pt,
  arc=2pt, left=10pt, right=10pt, top=7pt, bottom=7pt,
  before upper={\textcolor{bpshijau!60!black}{\sffamily\bfseries #1.}~}}

\newtcolorbox{kotakperhatian}[1][Perhatian]{
  breakable, enhanced, colback=bpsoranye!8, colframe=bpsoranye, boxrule=0pt, leftrule=4pt,
  arc=2pt, left=10pt, right=10pt, top=7pt, bottom=7pt,
  before upper={\textcolor{bpsoranye!80!black}{\sffamily\bfseries #1.}~}}

\definecolor{merahsalah}{HTML}{C0392B}
\newtcolorbox{kotaksalah}[1][Kesalahan Umum]{
  breakable, enhanced, colback=merahsalah!6, colframe=merahsalah, boxrule=0pt, leftrule=4pt,
  arc=2pt, left=10pt, right=10pt, top=7pt, bottom=7pt,
  before upper={\textcolor{merahsalah}{\sffamily\bfseries #1.}~}}

\newtcolorbox{kotaktips}[1][Tips dan Praktik Baik]{
  breakable, enhanced, colback=bpshijau!7, colframe=bpshijau, boxrule=0pt, leftrule=4pt,
  arc=2pt, left=10pt, right=10pt, top=7pt, bottom=7pt,
  before upper={\textcolor{bpshijau!55!black}{\sffamily\bfseries #1.}~}}

\newtcolorbox{kotakvariasi}[1][Variasi dan Alternatif]{
  breakable, enhanced, colback=bpsbirumuda!14, colframe=bpsbiru, boxrule=0pt, leftrule=4pt,
  arc=2pt, left=10pt, right=10pt, top=7pt, bottom=7pt,
  before upper={\textcolor{bpsbiru}{\sffamily\bfseries #1.}~}}

\newtcolorbox{kotakrumus}{
  breakable, enhanced, colback=bpsnavy!3, colframe=bpsnavy!35, boxrule=0.7pt,
  arc=3pt, left=10pt, right=10pt, top=6pt, bottom=6pt}

\newtcolorbox{kotaktugas}{
  breakable, enhanced, colback=bpsoranye!5, colframe=bpsnavy, boxrule=1pt, arc=3pt,
  left=10pt, right=10pt, top=8pt, bottom=8pt}

\newtcolorbox{kotakkunci}{
  breakable, enhanced, colback=bpshijau!6, colframe=bpshijau, boxrule=0.9pt, arc=3pt,
  left=10pt, right=10pt, top=9pt, bottom=8pt,
  title={\sffamily\bfseries Kunci \& Pembahasan}, coltitle=white, colbacktitle=bpshijau,
  fonttitle=\small, attach boxed title to top left={xshift=8pt, yshift=-3pt},
  boxed title style={boxrule=0pt, arc=2pt}}

% ---------- Kotak & perintah Daftar Istilah ----------
\newtcolorbox{kotakistilah}{
  breakable, enhanced, colback=bpsnavy!4, colframe=bpsnavy, boxrule=0pt,
  leftrule=4pt, arc=2pt, left=10pt, right=10pt, top=8pt, bottom=8pt}
% \istilahitem{istilah}{inggris}{definisi}
\newcommand{\istilahitem}[3]{%
  \noindent\textbf{\color{bpsnavy}#1}%
  \ifx\relax#2\relax\else\ {\small\color{abuteks}(#2)}\fi\ \textemdash\ #3\par
  \vspace{3pt}}
% \istilahkbbi{catatan KBBI}
\newcommand{\istilahkbbi}[1]{%
  {\small\color{bpsoranye}\textbf{KBBI:}}\ {\small\color{abuteks}#1}\par\vspace{5pt}}

% ---------- Penomoran Contoh, Kode, Latihan, Tugas ----------
\newcounter{contohctr}
\newcommand{\contoh}[1]{%
  \refstepcounter{contohctr}%
  \par\vspace{12pt}%
  \noindent\begin{tcolorbox}[enhanced, colback=bpsoranye!12, colframe=bpsoranye,
    boxrule=0pt, leftrule=5pt, arc=2pt, left=9pt, right=9pt, top=4pt, bottom=4pt]
  \sffamily\bfseries\color{bpsnavy}Contoh \arabic{section}.\arabic{contohctr}%
  \textcolor{bpsoranye}{\ \ $\vert$\ \ }#1
  \end{tcolorbox}\par\vspace{2pt}\nopagebreak}

\newcounter{kodectr}
\newcommand{\kode}[2][R]{%
  \refstepcounter{kodectr}%
  \par\vspace{8pt}\nopagebreak
  \def\lblbahasa{#1}%
  \ifx\lblbahasa\bahasaPython
    \noindent\colorbox{pykuning}{\textcolor{white}{\sffamily\bfseries\small\;Kode \arabic{section}.\arabic{kodectr} · Python\;}}%
  \else
    \noindent\colorbox{bpsbiru}{\textcolor{white}{\sffamily\bfseries\small\;Kode \arabic{section}.\arabic{kodectr} · R\;}}%
  \fi
  \;{\small\bfseries\color{bpsnavy}#2}%
  \par\vspace{-3pt}\nopagebreak}
\def\bahasaPython{Python}

\newcounter{latihanctr}
\newcommand{\latihan}[1]{%
  \refstepcounter{latihanctr}%
  \par\vspace{12pt}%
  \noindent\begin{tcolorbox}[enhanced, colback=bpshijau!12, colframe=bpshijau,
    boxrule=0pt, leftrule=5pt, arc=2pt, left=9pt, right=9pt, top=4pt, bottom=4pt]
  \sffamily\bfseries\color{bpsnavy}Latihan \arabic{section}.\arabic{latihanctr}%
  \textcolor{bpshijau!60!black}{\ \ $\vert$\ \ }#1
  \end{tcolorbox}\par\vspace{2pt}\nopagebreak}

\newcounter{tugasctr}
\newcommand{\tugas}[1]{%
  \refstepcounter{tugasctr}%
  \par\vspace{10pt}%
  \noindent\begin{tcolorbox}[enhanced, colback=bpsnavy!8, colframe=bpsnavy,
    boxrule=0pt, leftrule=5pt, arc=2pt, left=9pt, right=9pt, top=4pt, bottom=4pt]
  \sffamily\bfseries\color{bpsnavy}Tugas \nomormodul.\arabic{tugasctr}%
  \textcolor{bpsnavy}{\ \ $\vert$\ \ }#1
  \end{tcolorbox}\par\vspace{2pt}\nopagebreak}

% pastikan environment Shaded ada walau dokumen tak punya blok kode
\makeatletter
\@ifundefined{Shaded}{\newenvironment{Shaded}{}{}}{}
\@ifundefined{Highlighting}{\newenvironment{Highlighting}[1][]{}{}}{}
\makeatother
% ---------- Gaya blok kode R (environment Shaded dari pandoc) ----------
\renewenvironment{Shaded}{%
  \begin{mdframed}[backgroundcolor=codebg, linecolor=bpsbiru, linewidth=3pt,
    topline=false, bottomline=false, rightline=false,
    innerleftmargin=10pt, innerrightmargin=8pt, innertopmargin=7pt,
    innerbottommargin=7pt, skipabove=4pt, skipbelow=6pt]%
  \small}{\end{mdframed}}

% ---------- Gaya blok output R ----------
\newenvironment{routbox}{%
  \begin{mdframed}[backgroundcolor=outbg, linecolor=bpshijau, linewidth=3pt,
    topline=false, bottomline=false, rightline=false,
    innerleftmargin=10pt, innerrightmargin=8pt, innertopmargin=6pt,
    innerbottommargin=6pt, skipabove=2pt, skipbelow=8pt]%
  \footnotesize}{\end{mdframed}}

% ---------- Daftar berpoin ----------
\setlist[itemize]{leftmargin=1.4em, itemsep=2pt, topsep=3pt}
\setlist[enumerate]{leftmargin=1.7em, itemsep=2pt, topsep=3pt}

% ---------- Hyperref (dimuat pandoc setelah preamble) ----------
\AtBeginDocument{\hypersetup{colorlinks=true, linkcolor=bpsbiru, urlcolor=bpsbiru, citecolor=bpsbiru}}

% ---------- Halaman cover pembatas ----------
% \coverpembatas[BADGE]{nomor}{judul}{pertemuan}{sub-cpmk}
\newcommand{\coverpembatas}[5][]{%
  \begin{titlepage}
  \thispagestyle{empty}
  \begin{tikzpicture}[remember picture, overlay]
    \fill[white] (current page.south west) rectangle (current page.north east);
    % pita diagonal atas
    \fill[bpshijau!75!black] (current page.north west) -- ++(0,-7.4cm) --
      ($(current page.north east)+(0,-4.4cm)$) -- (current page.north east) -- cycle;
    \fill[bpshijau, opacity=0.9] (current page.north west) -- ++(0,-5.5cm) --
      ($(current page.north east)+(0,-2.7cm)$) -- (current page.north east) -- cycle;
    % ornamen
    \fill[bpsoranye, opacity=0.13] ($(current page.east)+(-1.6cm,1.2cm)$) circle (3.9cm);
    \fill[bpshijau, opacity=0.10] ($(current page.south east)+(-4.6cm,3.0cm)$) circle (2.6cm);
    % pita bawah
    \fill[bpsnavy] (current page.south west) rectangle ($(current page.south east)+(0,1.6cm)$);
    \fill[bpsoranye] ($(current page.south west)+(0,1.6cm)$) rectangle ($(current page.south east)+(0,1.78cm)$);
    % tiga blok warna BPS
    \fill[bpsbiru]   ($(current page.south west)+(2.5cm,2.7cm)$) rectangle ++(2.2cm,0.26cm);
    \fill[bpsoranye] ($(current page.south west)+(4.9cm,2.7cm)$) rectangle ++(2.2cm,0.26cm);
    \fill[bpshijau]  ($(current page.south west)+(7.3cm,2.7cm)$) rectangle ++(2.2cm,0.26cm);
    % identitas institusi
    \node[anchor=north west, text=white, align=left]
      at ($(current page.north west)+(2.5cm,-1.7cm)$)
      {\sffamily\Large POLITEKNIK STATISTIKA STIS\\[3pt]
       \sffamily\normalsize Program Studi D-IV Komputasi Statistik};
    \node[anchor=north west, text=white, align=left]
      at ($(current page.north west)+(2.5cm,-3.7cm)$)
      {\sffamily\small MODUL PRAKTIKUM \ \textbullet\ \ PEMROGRAMAN STATISTIKA (K242310)};
    % blok nomor modul
    \fill[bpsoranye] ($(current page.west)+(2.5cm,2.0cm)$) rectangle ++(0.18cm,3.7cm);
    \node[anchor=north west, text=bpsnavy] at ($(current page.west)+(3.0cm,5.75cm)$)
      {\sffamily\bfseries\fontsize{20}{22}\selectfont MODUL};
    \node[anchor=north west, text=bpsnavy] at ($(current page.west)+(2.9cm,4.85cm)$)
      {\sffamily\bfseries\fontsize{78}{78}\selectfont #2};
    % judul modul
    \node[anchor=north west, text=bpsnavy, align=left, text width=14.5cm]
      at ($(current page.west)+(2.5cm,1.5cm)$)
      {\sffamily\bfseries\fontsize{32}{38}\selectfont #3};
    % badge opsional
    \node[anchor=north west] at ($(current page.west)+(2.5cm,-1.1cm)$)
      {\ifx\relax#1\relax\else\colorbox{bpshijau}{\textcolor{white}{\sffamily\bfseries\ #1\ }}\fi};
    % keterangan
    \node[anchor=north west, text=abuteks, align=left, text width=13.5cm]
      at ($(current page.west)+(2.5cm,-2.1cm)$)
      {\sffamily\large\bfseries #4\\[7pt]
       \sffamily\normalsize\justifying #5};
    % penyusun
    \node[anchor=west, text=abuteks, align=left]
      at ($(current page.south west)+(2.5cm,3.55cm)$)
      {\sffamily\footnotesize\textbf{Penyusun:}\ \ Achmad Syahrul Choir, Rudi Salam, Christiana Anggraeni Putri};
    % teks pita bawah
    \node[anchor=west, text=white] at ($(current page.south west)+(2.5cm,0.8cm)$)
      {\sffamily\small Semester 3 (Ganjil)};
  \end{tikzpicture}
  \end{titlepage}%
  % penomoran berlanjut lintas modul (buku gabungan)
  \setcounter{page}{2}}

% --- penyesuaian khusus buku pedoman ---
\renewcommand{\sectionmark}[1]{\markboth{#1}{}}
\makeatletter
\@addtoreset{contohctr}{section}
\@addtoreset{kodectr}{section}
\@addtoreset{latihanctr}{section}
\makeatother
% kotak 'langkah' untuk instruksi bertahap
\newtcolorbox{kotaklangkah}[1][Langkah]{breakable, enhanced,
  colback=bpshijau!5, colframe=bpshijau, boxrule=0pt, leftrule=4pt, arc=2pt,
  left=10pt, right=10pt, top=7pt, bottom=7pt,
  before upper={\textcolor{bpshijau!70!black}{\sffamily\bfseries #1.}~}}
% Bagian II dimulai dari Bab 4



% penomoran berlanjut pedoman
\setcounter{page}{55}
